\originalchapter{22.02 作业二:波包与对易}
题源: Spring 2012 Problem Set 2, 2012年2月22日. 题 1、题 2(a)、题 5 部分含官方示范, 其余为编者独立解答.
\begin{exercise}\label{ex:2202-2-1}
原题 1 (已解). 求 $\int_{-\infty}^{\infty}(2\pi\sigma^2)^{-1/2}\e^{-x^2/(2\sigma^2)}\dd x$.
\end{exercise}
\begin{solution}
官方示范重编. 令 $I=\int\e^{-x^2}\dd x>0$. 平方后在二维极坐标积分, $I^2=\int_0^{2\pi}\dd\theta\int_0^\infty r\e^{-r^2}\dd r=\pi$, 因而 $I=\sqrt\pi$. 换元 $x=\sqrt2\sigma u$ 得题中积分为 1, 要求 $\sigma>0$.
\end{solution}
\begin{exercise}\label{ex:2202-2-2}
原题 2. 求 (a) $\sin(ax)$ (已解), (b) 上题高斯密度, (c) 矩形 $r(x)=1/(2a)$ ($|x|\le a$), 区间外零的 Fourier 变换. 原题最后一小项误重复写 b.
\end{exercise}
\begin{solution}
(a) 官方示范按本稿 Fourier 号复算为 $\sqrt{\pi/2}[\delta(k-a)-\delta(k+a)]/\ii$; 原示范的 $\ii$ 因子方向应与 $\e^{-\ii kx}$ 约定核对. (b) 配方或对 k 的积分求微分方程, 得 $\widetilde g(k)=\e^{-\sigma^2k^2/2}/\sqrt{2\pi}$. (c) 直接积分 $\int_{-a}^a\e^{-\ii kx}\dd x=2\sin ka/k$, 得 $\widetilde r=\sin(ka)/(ak\sqrt{2\pi})$, $k=0$ 按极限为 $1/\sqrt{2\pi}$.
\end{solution}
\begin{exercise}\label{ex:2202-2-3}
原题 3. 电子一维态 $\psi(x)=\pi^{-1/2}\cos(x/\sqrt2)$, 测动能可能得到什么值, 各概率是多少?
\end{exercise}
\begin{solution}
$-\psi''=\psi/2$, 故动能唯一值 $\hbar^2/(4m_e)$, 概率 1. 这里 x、波数按题设单位理解. 全实轴上 cos 非可归一化, 该陈述采用广义态或适当周期盒, 不以 $\pi^{-1/2}$ 宣称在全空间归一.
\end{solution}
\begin{exercise}\label{ex:2202-2-4}
原题 4. 中子态 $\psi=\pi^{-1/2}\sin(x/\sqrt3)$, (a) 写动能基表示; (b) 给动能测量结果和概率.
\end{exercise}
\begin{solution}
以 $k=1/\sqrt3$ 的两动量广义态, sin 的两系数相位相反, 模平方各 $1/2$. 两者动能同为 $\hbar^2/(6m_n)$, 所以能量测量该值的总概率是 1. 动量简并不应在列能量结果时重复计数.
\end{solution}
\begin{exercise}\label{ex:2202-2-5}
原题 5. $\psi=(2\pi a^2)^{-1/4}\e^{-x^2/(4a^2)}\e^{-\ii k_0x}$. 求 (a) $\langle x\rangle$, (b) $\Delta x^2$, (c) 动量表象与结果概率, (d) 用动量概率积分及直接算符作用两法求 $\langle p\rangle$, (e) 补齐 $\langle p^2\rangle$ 与 $\Delta p^2$; (f) 给不确定度乘积及 $a\to0$ 行为. 原题字母编号有重复.
\end{exercise}
\begin{solution}
(a),(b) 为 $0,a^2$. (c) $\widetilde\psi(k)=(2a^2/\pi)^{1/4}\e^{-a^2(k+k_0)^2}$, 动量密度 $P(p)=\sqrt{2a^2/\pi}\e^{-2a^2(p/\hbar+k_0)^2}/\hbar$. (d) 对称高斯的均值为 $-\hbar k_0$. 直接用 $p\psi=(\ii\hbar x/(2a^2)-\hbar k_0)\psi$ 并积分, 奇项零, 相同结果. (e) 再作用一次或积分高斯二阶矩得 $\langle p^2\rangle=\hbar^2k_0^2+\hbar^2/(4a^2)$, 所以 $\Delta p^2=\hbar^2/(4a^2)$. (f) $\Delta x\Delta p=\hbar/2$, $a\to0$ 时动量展宽无界. 源题称 $p_0=+\hbar k_0$ 与实际负指数不一致, 本稿按给定波函数保留负均值.
\end{solution}
\begin{exercise}\label{ex:2202-2-6}
原题 6. 一维宇称 $Pf(x)=f(-x)$. (a) 对自由 $H=p^2/(2m)$ 证明 $[H,P]=0$. (b) 从 $\e^{\pm\ii kx}$ 构造共同本征函数并给本征值.
\end{exercise}
\begin{solution}
(a) $(Pf)''(x)=f''(-x)$, 所以 $HPf=PHf$. (b) cos 与 sin 分别具有 P 本征值 $+1,-1$, 两者动能同为 $\hbar^2k^2/(2m)$. $k=0$ 时 sin 恒零, 不能作为本征态.
\end{solution}
\begin{exercise}\label{ex:2202-2-7}
原题 7. 哪些动量分量同 H 对易: (a) $V=\e^{-x}(y^3+3z)$, (b) $V=a\sin(bx)/(bx)$, (c) $V=b\e^{-(x^2+y^2)/c^2}$? (d) 哪些势允许能量与三维宇称共同本征态? 参数为常数, $bx=0$ 按连续延拓.
\end{exercise}
\begin{solution}
$[p_i,H]=-\ii\hbar\partial_iV$. 一般非零参数下 (a) 无动量分量对易; (b) $p_y,p_z$; (c) $p_z$. 对 (b),(c), $V(-\mathbf r)=V(\mathbf r)$, 可以选共同宇称能量基; (a) 一般没有. 参数使势变常数或为零时会额外获得对称性, 必须另作退化情况.
\end{solution}
